\documentclass{tutodoc}

\usepackage{../preamble.cfg}

\usepackage{../main/core-ctxt-dsl}
\usepackage{../main/main}


% TESTING LOCAL IMPLEMENTATION %

\usepackage{data-1D}


\begin{document}

\section{Tableaux d'images de fonctions à une seule variable}

\subsection{Descriptions des contextes à utiliser}

La saisie de tableaux, forcément en mode mathématique, d'images de fonctions à une seule variable nécessitent de fournir les valeurs pivots qui sont celles dont on veut indiquer les images. Ceci peut se faire de deux façons.
%
\begin{enumerate}
	\item \tdoclatexin{xvals = x1 , x2 , ... , xn} indiquent les valeurs pour la variable $x$, qui est celle par défaut.

	\item \tdoclatexin{xvals = mavar : x1 , x2 , ... , xn} permet de renommer la variable en $mavar$.
\end{enumerate}


Pour les images, il existe deux modes de saisie, mais avec une restriction forte, car une seule série d'images peut utiliser le nom par défaut.
%
\begin{enumerate}
	\item \tdoclatexin{imgs = i1 , i2 , ... , in} indiquent les images pour la fonction $f(x)$, qui est celle par défaut.

	\item \tdoclatexin{imgs = mafonc : i1 , i2 , ... , in} indique les images pour la fonction $mafonc$.
\end{enumerate}

En résumé, tout se passe via \tdoclatexin{xvals = ...} pour les valeurs pivots, et \tdoclatexin{imgs = ...} pour les images.
%
Il est possible d'employer \tdoclatexin{x} et \tdoclatexin{i} comme raccourcis de \tdoclatexin{xvals} et \tdoclatexin{imgs}.


\end{document}
